\section{Conclusion}
\label{sec:conclusion}

Our experiments show that compact relational models can transfer across
databases: \model{} achieves performance comparable to \RT{} with 12.5\%
of its parameter count, and higher average accuracy at a similar model size.
The case study in Section~\ref{sec:case_study} shows that two-hop neighborhoods
recover most of the full-context performance, historical task labels are a
major source of predictive information, and disabling relation conditioning
reduces average accuracy.

These findings also reveal limitations and directions for future work.
Both \model{} and \RT{} deteriorate substantially when the query entity's
historical task labels are hidden. Improving prediction for entities with
little or no label history therefore remains a central challenge. Future
work could explore training with missing label histories to encourage
greater use of activity records and evidence from related entities.
Our experiments show gains from increasing model size, but pretraining
remains limited to the RelBench databases. An important direction is to
expand the scale and diversity of pretraining data and examine whether
these gains extend to a broader range of unseen schemas and tasks.
Finally, our evaluation is limited to binary entity classification;
extending the prediction mechanism and evaluating transfer on regression
and recommendation tasks remain open directions.

We introduced \model{}, combining transferable relation representations
from RelationGraph, RowGraph reasoning, and shared class-candidate
prediction. Together, these components provide a compact architecture for
entity classification across unseen relational databases, while our
analysis identifies reliance on entity label history as a key challenge
for further progress.
